%% kapa: exam=1 round=31 year=2020 period=1 form=A subject=real_estate subject_ko=부동산학원론 seq=16 number=96
%% source: Q-Net 감정평가사 기출문제 (cst00309) artlSeq=5207219 file=extracted/1008/5207219/A형/A형/2020년도 제31회 감정평가사 1차 시험_1교시_A형.pdf
%% conversion: kordoc 4.24.51 -> tools/gen1.py (auto); review=proofread-sonnet; flags=none
\begin{question}{96}
토지이용계획과 용도지역지구제에 관한 설명으로 옳지 않은 것은?
\begin{choices}
  \item 용도지역지구제는 토지이용규제의 대표적인 예로 들 수 있다.
  \item 용도지역지구제는 특정 토지를 용도지역이나 용도지구로 지정한 후 해당 토지의 이용을 지정목적에 맞게 제한하는 제도이다.
  \item 토지이용계획은 토지이용규제의 근간을 이루지만 법적 구속력을 가지고 있지는 않다.
  \item 용도지역지구제는 토지이용계획의 내용을 실현하는 수단으로서, 도시ㆍ군관리계획의 내용을 구성한다.
  \item 용도지역지구제에 따른 용도 지정 후, 관련법에 의해 사인의 토지이용이 제한되지 않는다.
\end{choices}
\end{question}
\answer{5}
